\section{結論}
\label{sec:conclusion}

本稿は、実会場の 3DGS 再構成から
テニスコート検出のための学習視点分布を作る手法 \Method を提示した。
中心にあるのはアーキテクチャではなく、
シーン単位で一度だけ検証されるメートル系アライメントである。
地上面へ持ち上げた線証拠から共通尺度の規制テンプレートを二面当てはめ、
fit と holdout を分離して凍結検証し、
トポロジーと意味セグメントの被覆を満たさなければシーンを棄却する。
この停止条件が、
誤った幾何教師が数千枚単位で生成されることを防ぐ。

受理された一つの変換は、
復元カメラの凸包で制限した三次元軌道からの
RGB、カメラ行列、シーン--コート変換、
幾何妥当性、描画支持、
物理的なキーポイント同一性を同時に支配する。
人手の費用は画像あたりからシーンあたりへ移り、
視点数と多様性に比例して償却される。

本研究の中心的教訓は、
ここで合成データが有用なのは
画像を似せて増やすからではなく、
各画像の背後にある幾何を
ほぼ無視できる限界費用で露わにするからだ、という点である。

ただし本稿は、
視点拡張が実写の非放送視点への汎化を改善することを
まだ示していない。
本稿が示したのは、
実写 held-out validation と合成 trajectory-group test で
同一のアライメント手続きを適用した交差モデル比較である。
実写 validation では \Method が僅かに上回り
（PCK@0.01 で \resRealOursPckOne 対 \resRealTcdPckOne、
対応誤差の中央値で \resRealOursPairMedian px 対 \resRealTcdPairMedian px）、
合成 test では差が大きく
（PCK@0.01 で \resSynthOursPckOne 対 \resSynthTcdPckOne）、
ベースラインの有効対応は \resSynthTcdPairN 点分にとどまった。
しかしこの二つのモデルはアーキテクチャも学習データも異なる。
したがってこの差は視点拡張の因果効果を分離しない。
checkpoint 選択時の合成 validation は別結果であり、
その run は完走していない。
同一条件の狭域・広域比較、
実写のみの同一アーキテクチャ比較、
未知会場の分離、複数 seed が
仮説を判定する次の段階である。
本稿の貢献は、
その比較を意味のある形で行えるデータ生成契約と評価手続きを、
失敗時の停止条件を含めて固定し、
現在の checkpoint で実際に測定して公開したことにある。
